% gt_part_M4.tex -- [GT] the M4 skeleton: G5 (the Hecke ring and its relations) and G9 (level N, nebentypus) headers.
% Slot 19, stretch package L19-5 (docs/implementation/slot19/L19-5_M4_skeleton.md).  Read by lean_census.py (gt_part_*.tex).
% Lean: lean/GL2/M4/Skeleton.lean (namespace GL2.M4).
% Slot 19 (L19-5): every row was then `planned` (no longer true; see below).  Each \lean{} name is a sorry-free `def ... : Prop` (the frozen statement) or a
% sorry-free definition, so the census resolves it and counts the row `planned`; it never counts `stated` and no `sorry`
% exists (L19-0 §5.4).  Slot 20 (L20-1, M4-a / M4-b): lem:m4-tauR-exists, lem:m4-cosetdata, prop:diamond-model and
% prop:fullassembly-model are `proved` (Lean: lean/GL2/M4/Data.lean); def:m4-hecke stayed `planned` until F3-20 ruled on Tfull
% (L20-6 audit §4: Tfull stands as the assembled-form definition); `proved` since the L20-6 repairs, so no row here is `planned`.
% Slot 20 (L20-2, M4-c/d/e, branch H): the G5 rows are `proved_mod_lit` over the Part-L rows lit:hecke-Tp-prim / -coprime /
% -commute (gt_part_L.tex, GL2/Literature/HeckeModel.lean; Shimura locators verified, R20-L); Lean: lean/GL2/M4/Relations.lean.
% Slot 20 (L20-5, M4-f, branch H): def:m4-classical (proved) and thm:classical-model (proved_mod_lit over lit:hecke-classical,
% GL2/Literature/HeckeClassical.lean; Brown locators verified, Ash-Stevens leg not as cited, R20-L); Lean: lean/GL2/M4/{ClassicalDefs,Classical}.lean.
% Slot 20 (L20-3, M4-g/h): prop:levelN-shapiro-model, def:levelellmodel-model, prop:levelcorrespondence-model, def:nebdata-model,
% prop:nebcorr-model, lem:m4-fullrec are `proved`, thm:nebrelations-model `proved_mod_lit`; Lean: lean/GL2/M4/LevelN.lean.
% Slot 20 (L20-4, M4-i): def:merelset-model, lem:merelpath-model, prop:heilbronn-model, chk:merel-n5 are `proved`;
% Lean: lean/GL2/M4/{Merel,Heilbronn}.lean, lean/GL2/Checkers/{MerelCheck,MerelData,MerelN5}.lean.
% Slot 20 (L20-6 repairs, docs/implementation/slot20/L20-6_repairs.md): the F3-20 findings MAJOR 1 (thm:classical-model is the
% weak form; rmk:m4-g5-headers lists [C] thm:classical as open) and MINOR 2-10 (def:m4-hecke's scope sentence and leanok,
% clause attributions, scope and choice-freeness caveats, the MEREL_DUMP provenance hash, cosmetics).
% Slot 20 (close delta, docs/implementation/slot20/close_delta.md): rmk:m4-g5-headers' obl:algebra item corrected after L20-H
% (the FLT harvest has HeckePair.HeckeAlgebra and HeckeCohomology.heckeH1); text only, no row, \lean{} name or status changes.
% The rows still `planned` after slot 20 are listed by lean_census.sh (status column); no row here is `stated`.
% Slot 20 proves a row by
% adding `theorem <name> : <name>_stmt ...`, adding <name> to the row's \lean{}, and setting \leanok on the statement and the
% proof (the row then counts `proved`; the `_stmt` definition stays linked as the statement the theorem matches).
% Labels: [GT]-only labels (suffix -model or m4:); each names the [C] v1.18.0 clause it states, in a "% states" line.

\chapter{Part M4. The Hecke ring and level \texorpdfstring{$N$}{N} on the model}\label{ch:m4}

\paragraph{Setting.} As in \ref{def:heckeP}: $G=\mathrm{GL}_2(\mathbb Q)$, $\Gamma=\mathrm{SL}_2(\mathbb Z)$ (\texttt{ΓQ}), $\Delta$ the
integral matrices (\texttt{ΔQ}), the seven-vertex model, $\tau$ over $R$ (\texttt{TauR}), and any coefficient module
$\rho\colon\Delta\to\mathrm{End}_R(V)$. $H^n$ below is $H^n(\mathrm{Hom}_\Gamma(C_\bullet(X_{\mathrm{bulk}}),V))$, which is
$H^n_{\mathrm{eq}}(K,\mathrm{Rel};L)$ by $\Psi$ (\ref{def:heckeP}, \texttt{heckeRel}) and $H^n(\mathrm{SL}_2(\mathbb Z);V)$ by
\ref{thm:model7-rmkorbifold} when $6\in R^\times$. The scalar $c$ is the one by which $\rho(pI)$ acts; for
$\rho=\rho_n=\mathrm{Sym}^n$ it is $p^n$ ([C] \texttt{prop:symconv}(2)), and then $p\,c=p^{n+1}$.

\section{The objects}

\begin{definition}[$T_{\Gamma g\Gamma}$ on $H^n$; the Hecke algebra; primitive and full $p$-power operators]\label{def:m4-hecke}
  % states [C] v1.18.0 def:operators (1)-(3) on the model; eq:fullop through eq:fullassembly (see the note)
  \lean{GL2.M4.HM, GL2.M4.Tdc, GL2.M4.heckeAlg, GL2.M4.IsScalarQ, GL2.M4.Tprim, GL2.M4.Tfull, GL2.M4.Tp, GL2.M4.Tp_eq,
    GL2.M4.TdcOf, GL2.M4.TdcOf_eq}
  \leanok
  \uses{def:heckeP, thm:welldef-b}
  $T_{\Gamma g\Gamma}=\texttt{Tdc}\,\rho\,cd\,\tau\,n$ is the map that \texttt{heckeTildeR} induces on $H^n$. \texttt{heckeAlg} is the
  $R$-subalgebra of $\mathrm{End}_R(H^n)$ generated by every $T_{\Gamma g\Gamma}$, $g\in\Delta$, every coset datum, every $\tau$.
  \texttt{Tprim} is $T_{\Delta^{\mathrm{prim}}_{p^\mu}}$ on \texttt{cdP} (and $1$ at $\mu=0$). \texttt{Tfull} is
  $T_{p^\nu}=\sum_{0\le2j\le\nu}c^jT_{\Delta^{\mathrm{prim}}_{p^{\nu-2j}}}$.
  \emph{Scope} (F3-20 \S4, the ruling on \texttt{Tfull}). \texttt{Tfull} is \emph{defined} by \texttt{eq:fullassembly} (with
  $T(p,p)$ acting by $c$); its agreement with [C] \texttt{eq:fullop} (the sum over all double cosets of determinant $p^\nu$) uses
  the elementary-divisor enumeration of those double cosets (they are exactly $\Gamma p^j\mathrm{diag}(1,p^{\nu-2j})\Gamma$,
  $0\le2j\le\nu$, pairwise distinct); every row below is stated for \texttt{Tfull}. Each summand is
  \ref{prop:fullassembly-model}. [C] \texttt{eq:fullop} receives no \texttt{\textbackslash LEAN} marker. \emph{Slot 22 (M5):} the
  enumeration is \ref{thm:m5-dc-enum}, and $\sum_iT_{\Gamma g_i\Gamma}=\texttt{Tfull}$ over every complete irredundant system of
  representatives is \ref{thm:m5-fullop}.
  \texttt{Tp} ($6\in R^\times$) is $T_{\Gamma\mathrm{diag}(1,p)\Gamma}$ on \texttt{cdP}~$p$~$1$ with a chosen $\tau$
  (\ref{lem:m4-tauR-exists}), and \texttt{TdcOf}~$g$ is $T_{\Gamma g\Gamma}$ with a chosen coset datum (\ref{lem:m4-cosetdata}) and a
  chosen $\tau$: the choice-free operators. \texttt{Tp\_eq} and \texttt{TdcOf\_eq} (slot 20, L20-1) say that they equal
  \texttt{Tdc} at \emph{every} coset datum and every $\tau$ (\ref{thm:welldef-b} and \texttt{heckeTildeR\_indep\_cd}).
\end{definition}
\begin{proof}
  \leanok
  Definitions (sorry-free); \texttt{Tp\_eq} and \texttt{TdcOf\_eq} are proved in \texttt{GL2/M4/Data.lean} (L20-1).
\end{proof}

\section{Existence of the data}

\begin{lemma}[$\tau$ over $R$ exists for every $g$]\label{lem:m4-tauR-exists}
  % states [C] v1.18.0 lem:tauexists (2), existence over R with 6 a unit, in the form heckeTildeR reads
  \lean{GL2.M4.tauR_nonempty_stmt, GL2.M4.tauR_nonempty}
  \leanok
  \uses{lem:tauexists-2-model, def:heckeP}
  For $6\in R^\times$ and every $g\in\mathrm{GL}_2(\mathbb Q)$ a \texttt{TauR}~$R$~$g$ exists. (Route: \texttt{model7\_tauexists\_2} on
  \texttt{fIncl}~$g$, \texttt{fConj}~$g$ gives a chain map of the restricted resolutions over the augmentation; its components
  are a \texttt{TauR}.)
\end{lemma}
\begin{proof}
  \leanok
  \ref{lem:tauexists-2-model} at $f_1=\texttt{fIncl}~g$, $f_2=\texttt{fConj}~g$ gives $\tau\colon P_1\to P_2$ with
  $\tau\circ\pi_2=\pi_1$. Its components are the maps of a \texttt{TauR}; the chain-map square in degree $n+1\to n$ is (T2), and the
  degree-$0$ component of $\tau\circ\pi_2=\pi_1$, evaluated on elements, is $\varepsilon\circ\tau_0=\varepsilon$ (T3).
\end{proof}

\begin{lemma}[Coset data exist for every $g\in\Delta$]\label{lem:m4-cosetdata}
  % states [C] v1.18.0 def:hecke (the coset decomposition), prop:cosets, def:cosetrepsN
  \lean{GL2.M4.cosetData_nonempty_stmt, GL2.M4.cosetData_nonempty, GL2.M4.Γpp_finiteIndex}
  \leanok
  \uses{def:heckeP, prop:cosets-4}
  For every integral $g\in\mathrm{GL}_2(\mathbb Q)$, $\Gamma''=\Gamma\cap g\Gamma g^{-1}$ has finite index in $\Gamma$, and a
  transversal $s_1,\dots,s_m$ with $s_j^{-1}g\in\Delta$ exists (a \texttt{CosetData}).
\end{lemma}
\begin{proof}
  \leanok
  (L20-0 §4, the mathlib route; not the Smith form, not \texttt{def:cosetrepsN}.) mathlib's \texttt{exists\_Gamma\_le\_conj} at
  $g^{-1}$ and $M=1$ gives $N\neq0$ with $g^{-1}\gamma g\in\mathrm{SL}_2(\mathbb Z)$ for every $\gamma\in\Gamma(N)$, so
  $\Gamma(N)\le\Gamma''$ and $\Gamma''$ has finite index (\texttt{$\Gamma$pp\_finiteIndex}). Enumerate the left cosets $\Gamma/\Gamma''$ by
  $\mathrm{Fin}\,m$ and let $s_i$ be the inverse of a representative of the $i$-th: then $\gamma\in\Gamma''s_i$ iff
  $\gamma^{-1}\Gamma''$ is the $i$-th coset, which gives existence and uniqueness. $\delta_i=s_i^{-1}g\in\Delta$ because
  $\Gamma\subseteq\Delta$ and $\Delta$ is a monoid. The datum is not constructive; nothing depends on it (\texttt{TdcOf\_eq}).
\end{proof}

\section{G5: the relations}

\begin{proposition}[Central operator on the model]\label{prop:diamond-model}
  % states [C] v1.18.0 prop:diamond (Delta of N I acts by rho(N I))
  \lean{GL2.M4.diamond_stmt, GL2.M4.diamond, GL2.M4.diamond_id, GL2.M4.tauScalar, GL2.M4.central_of_scalar,
    GL2.M4.cdOne, GL2.M4.tauOne, GL2.M4.heckeTildeR_one, GL2.M4.HmapOf_one}
  \leanok
  \uses{def:m4-hecke, prop:fullassembly-model}
  If $g=N\cdot I$ and $\rho(g)=c$, then $T_{\Gamma g\Gamma}=c$ on $H^n$ for every coset datum and every $\tau$, $6\in R^\times$.
  \emph{Restatement.} [C] needs no hypothesis on $\ell$, because it uses $\tau=\mathrm{id}$. For an arbitrary $\tau$ the route
  passes through \ref{thm:welldef-b}, hence $6\in R^\times$. The $\tau=\mathrm{id}$ form, without that hypothesis, is
  \texttt{diamond\_id} ([C]'s scope): there $\tau$ is \texttt{tauScalar}, the identity maps (for $g$ central,
  $g^{-1}hg=h$), and every coset datum is allowed.
\end{proposition}
\begin{proof}
  \leanok
  $g=g\cdot1$ with $g$ central. By \texttt{heckeTildeR\_indep\_cd} the coset datum may be \texttt{cdScale} of \texttt{cdOne} (one
  coset, $s=1$); by the scaling identity (\texttt{heckeTildeR\_scale}, see \ref{prop:fullassembly-model}) the operator is
  $c\cdot\widetilde T_{\Gamma1\Gamma}$ with $\tau=\mathrm{id}$, and $\widetilde T_{\Gamma1\Gamma}=\mathrm{id}$ on cochains of degree
  $n\ge0$ (\texttt{heckeTildeR\_one}). This is \texttt{diamond\_id}, with no hypothesis on $R$. For an arbitrary $\tau$,
  \ref{thm:welldef-b} ($6\in R^\times$) replaces $\tau$ by \texttt{tauScalar}.
\end{proof}

\begin{proposition}[Primitive-to-full assembly on the model, the summand]\label{prop:fullassembly-model}
  % states [C] v1.18.0 prop:fullassembly: T_{p^j I . Delta^prim_{p^mu}} = p^{jn} T_{Delta^prim_{p^mu}}
  \lean{GL2.M4.fullassembly_stmt, GL2.M4.fullassembly, GL2.M4.transport_indep_cd, GL2.M4.heckeTildeR_indep_cd,
    GL2.M4.cdScale, GL2.M4.tauScale, GL2.M4.heckeTildeR_scale, GL2.M4.HmapOf_smul_of}
  \leanok
  \uses{def:m4-hecke, thm:welldef-b}
  For $g=p^j\,\mathrm{diag}(1,p^\mu)$, $\mu\ge1$, any coset datum and any $\tau$ for $g$, any $\tau'$ for
  $\mathrm{diag}(1,p^\mu)$, and $6\in R^\times$: $T_{\Gamma g\Gamma}=c^j\,T_{\Delta^{\mathrm{prim}}_{p^\mu}}$ on $H^n$.
  At cochain level with $\tau=\texttt{tauScale}~\tau'$ the identity holds with no hypothesis on $R$
  (\texttt{heckeTildeR\_scale}), which is [C]'s scope; $6\in R^\times$ enters only because the statement allows an arbitrary
  $\tau$ (\ref{thm:welldef-b}).
\end{proposition}
\begin{proof}
  \leanok
  \emph{Scaling identity.} Let $z\in\Delta$ be central with $\rho(z)=c$, and $g=zg'$. Then $g^{-1}hg=g'^{-1}hg'$, so
  $\Gamma''(g)=\Gamma''(g')$ and a coset datum of $g'$ is one of $g$ with $\delta_j=z\delta'_j$ (\texttt{cdScale}), and a $\tau$ for
  $g'$ is one for $g$ (\texttt{tauScale}, the same maps). With these, $\rho(\delta_j)=c\,\rho(\delta'_j)$ gives
  $\widetilde T_g=c\,\widetilde T_{g'}$ at cochain level (\texttt{heckeTildeR\_scale}), hence on $H^n$
  (\texttt{HmapOf\_smul\_of}).
  \emph{Here} $z=(pI)^j$ (central, $\rho(z)=c^j$) and $g'=\mathrm{diag}(1,p^\mu)$, so $g=zg'$. The given coset datum is replaced by
  \texttt{cdScale} of \texttt{cdP}~$p$~$\mu$ at cochain level (\texttt{heckeTildeR\_indep\_cd}: [C] \texttt{thm:welldef}(a) for
  $\tau$ over $R$, no hypothesis on $R$, by \texttt{CosetData.reindex}), and the given $\tau$ by \texttt{tauScale}~$\tau'$ on
  cohomology (\ref{thm:welldef-b}, $6\in R^\times$).
\end{proof}

\begin{theorem}[$T_p^2=T_{\Delta^{\mathrm{prim}}_{p^2}}+(p+1)c$, the clause $\mu=0$ of \texttt{hTp}]\label{thm:relations-psq-model}
  % states [C] v1.18.0 thm:relations (1) at nu = 1 with its caveat; obl:algebra (ii), (iii)
  \lean{GL2.M4.Tp_sq_stmt, GL2.M4.Tp_sq}
  \leanok
  \uses{def:m4-hecke, bridge:fullprimitive, lit:hecke-Tp-prim, lem:m4-relations-aux}
  For $6\in R^\times$: $T_p^2=T_{\Delta^{\mathrm{prim}}_{p^2}}+(p+1)c$ on $H^n$. This row is \texttt{hTp}'s clause $\mu=0$ for
  the frozen statement; the $p^nI$ difference is \emph{derived} in \ref{lem:m4-hTp-hrel} (F3-20 MINOR 5), and a
  \texttt{\textbackslash LEANlit} marker on [C] \texttt{thm:relations}(1)'s caveat or on \texttt{obl:algebra}(ii), (iii) cites that
  row. With $c=p^n$ this gives
  $T_{\Delta^{\mathrm{prim}}_{p^2}}=T_p^2-p^{n+1}-p^n$ (\texttt{obl:algebra}(iii)) and, for the full
  $T_{p^2}=T_{\Delta^{\mathrm{prim}}_{p^2}}+p^n$, $T_{p^2}=T_p^2-p^{n+1}$ (\texttt{obl:algebra}(ii)). Substituting either operator
  into the other's identity is off by exactly $p^nI$. \ref{bridge:fullprimitive} has this from the Hecke-ring identity as
  hypothesis \texttt{hrel}. This row keeps it as a named literature hypothesis (\textsf{proved\_mod\_lit}, the [LP]
  \textbf{H} branch; slot 20, L20-2).
  \gtmodlit{the transported product of primitive double cosets (Shimura Thm.~3.24(5) with Prop.~3.17, verified in R20-L), \ref{lit:hecke-Tp-prim}, as the hypothesis \gtlithyp{hTp} (\texttt{HeckeTpPrim})}
\end{theorem}
\begin{proof}
  \leanok
  The clause $\mu=0$ of \texttt{hTp} at the family $\tau_\mu=$ a chosen \texttt{TauR} (\ref{lem:m4-tauR-exists}); the given
  $\tau_1,\tau_2$ are replaced by the chosen ones by \texttt{heckeTildeR\_indep\_tau} (\ref{thm:welldef-b}).
\end{proof}

\begin{lemma}[Auxiliaries of the G5 rows (no literature)]\label{lem:m4-relations-aux}
  \lean{GL2.M4.mem_ΔQ_of_cd, GL2.M4.Tdc_eq_TdcOf, GL2.M4.Tprim_indep, GL2.M4.Tprim_pos, GL2.M4.Tfull_add_two,
    GL2.M4.Tfull_zero, GL2.M4.Tfull_one, GL2.M4.full_recursion_alg, GL2.M4.commute_of_mem_adjoin}
  \leanok
  \uses{def:m4-hecke, lem:m4-tauR-exists, lem:m4-cosetdata}
  (1) A coset datum for $g$ forces $g\in\Delta$, and $\texttt{Tdc}=\texttt{TdcOf}$. (2) \texttt{Tprim} does not depend on the family
  of $\tau$s ($6\in R^\times$). (3) The assembly $T_{p^{\nu+2}}=T_{\Delta^{\mathrm{prim}}_{p^{\nu+2}}}+c\,T_{p^\nu}$,
  $T_{p^0}=1$, $T_{p^1}=T_{\Delta^{\mathrm{prim}}_p}$. (4) In any $R$-algebra, from (3) and two identities of the shape of
  \ref{lit:hecke-Tp-prim} (as hypotheses of an algebra lemma): $P_1F_{\nu+1}=F_{\nu+2}+p\,c\,F_\nu$. (5) An $R$-subalgebra
  generated by pairwise commuting elements is commutative.
\end{lemma}
\begin{proof}
  \leanok
  (1) $g=s_i\cdot s_i^{-1}g$ for the coset of $1$; \texttt{TdcOf\_eq}. (2) \texttt{heckeTildeR\_indep\_tau}. (3) The definition of
  \texttt{Tfull}, \texttt{Finset.sum\_range\_succ'}. (4) Strong induction on $\nu$: $\nu=0$ is the clause $\mu=0$; $\nu=1$ the clause
  $\mu=1$; $\nu+2$: $P_1F_{\nu+3}=P_1P_{\nu+3}+c\,P_1F_{\nu+1}=P_{\nu+4}+pc\,P_{\nu+2}+c(F_{\nu+2}+pc\,F_\nu)
  =F_{\nu+4}+pc\,F_{\nu+2}$. (5) \texttt{Algebra.adjoin\_induction}, twice.
\end{proof}

\begin{lemma}[\texttt{hTp} gives the \texttt{hrel} of \texttt{bridge:fullprimitive}]\label{lem:m4-hTp-hrel}
  % L20-0 §3.2 check (a): the two literature routes (hrel of bridge:fullprimitive, hTp) agree on the model
  \lean{GL2.M4.hTp_imp_hrel_stmt, GL2.M4.hTp_imp_hrel, GL2.M4.hTp_prim_recursion_stmt, GL2.M4.hTp_prim_recursion}
  \leanok
  \uses{bridge:fullprimitive, lit:hecke-Tp-prim, lem:m4-relations-aux}
  At $A=\mathrm{End}_R(H^n)$, with $T_p=\mathtt{Tprim}_1$, $T_{p^2}=\mathtt{Tfull}_2$, $T_{\mathrm{prim}}=\mathtt{Tprim}_2$,
  $T(p,p)=c$: \texttt{hTp} gives \texttt{hrel} ($T_p^2=T_{p^2}+p\,T(p,p)$) and \texttt{hfull} ($T_{p^2}=T_{\mathrm{prim}}+T(p,p)$) of
  \ref{bridge:fullprimitive}; at $c=p^n$, \texttt{GL2.prim\_recursion} then gives $T_{p^2}=T_p^2-p^{n+1}$ and
  $T_{\mathrm{prim}}=T_p^2-p^{n+1}-p^n$ on the model. So the two literature routes agree.
  \gtmodlit{\ref{lit:hecke-Tp-prim} as the hypothesis \gtlithyp{hTp} (\texttt{HeckeTpPrim})}
\end{lemma}
\begin{proof}
  \leanok
  \texttt{hfull} is the assembly at $\nu=0$; \texttt{hrel} is the clause $\mu=0$ of \texttt{hTp} rewritten by \texttt{hfull}
  ($(p+1)c=c+p\,c$).
\end{proof}

\begin{theorem}[The recursion for the full operators]\label{thm:relations-1-model}
  % states [C] v1.18.0 thm:relations (1): T_{p^{nu+1}} = T_p T_{p^nu} - p^{n+1} T_{p^{nu-1}}
  \lean{GL2.M4.relations_stmt, GL2.M4.relations}
  \leanok
  \uses{def:m4-hecke, thm:relations-psq-model, prop:fullassembly-model, lit:hecke-Tp-prim, lem:m4-relations-aux}
  For $6\in R^\times$ and every family of $\tau$s: $T_{p^{\nu+2}}=T_p\,T_{p^{\nu+1}}-p\,c\,T_{p^\nu}$ on $H^n$, $\nu\ge0$.
  \gtmodlit{\ref{lit:hecke-Tp-prim} (Shimura Thm.~3.24(5), Prop.~3.17, verified in R20-L) as the hypothesis
  \gtlithyp{hTp} (\texttt{HeckeTpPrim})}
\end{theorem}
\begin{proof}
  \leanok
  Our own argument (the full / primitive separation): \ref{lem:m4-relations-aux}(4) with $P=\mathtt{Tprim}$,
  $F=\mathtt{Tfull}$, the assembly (3), and the two clauses of \texttt{hTp}.
\end{proof}

\begin{theorem}[Coprime multiplicativity, primitive form]\label{thm:relations-2-model}
  % states [C] v1.18.0 thm:relations (2) for the primitive double cosets of coprime M, N
  \lean{GL2.M4.coprime_stmt, GL2.M4.coprime}
  \leanok
  \uses{def:m4-hecke, lit:hecke-coprime, lem:m4-relations-aux}
  For $6\in R^\times$, coprime $M,N>0$, and any coset data and $\tau$s:
  $T_{\Gamma\mathrm{diag}(1,M)\Gamma}\,T_{\Gamma\mathrm{diag}(1,N)\Gamma}=T_{\Gamma\mathrm{diag}(1,MN)\Gamma}$ on $H^n$.
  \gtmodlit{\ref{lit:hecke-coprime} (Shimura Prop.~3.16, verified in R20-L) as the
  hypothesis \gtlithyp{hcop} (\texttt{HeckeCoprime})}
\end{theorem}
\begin{proof}
  \leanok
  Each \texttt{Tdc} is \texttt{TdcOf} (\ref{lem:m4-relations-aux}(1)); then \texttt{hcop}.
\end{proof}

\begin{theorem}[The operators commute]\label{thm:relations-comm-model}
  % states [C] v1.18.0 obl:algebra ("whence the {T_N} commute", all double cosets); thm:relations (1) first clause and
  % obl:algebra (i) are instances (F3-20 MINOR 4)
  \lean{GL2.M4.Tdc_commute_stmt, GL2.M4.Tdc_commute}
  \leanok
  \uses{def:m4-hecke, lit:hecke-commute, lem:m4-relations-aux}
  [C] \texttt{obl:algebra}'s consequence ``whence the $\{T_N\}$ commute'', for all double cosets (the first clause of [C]
  \texttt{thm:relations}(1), $T_pT_{p'}=T_{p'}T_p$, is an instance): for $6\in R^\times$, any $g,g'\in\Delta$ (here $\Delta$
  includes $\det<0$) with coset data and $\tau$s: $T_{\Gamma g\Gamma}T_{\Gamma g'\Gamma}=T_{\Gamma g'\Gamma}T_{\Gamma g\Gamma}$ on $H^n$.
  \gtmodlit{\ref{lit:hecke-commute} (Shimura Prop.~3.8, verified in R20-L)
  as the hypothesis \gtlithyp{hcomm} (\texttt{HeckeCommute})}
\end{theorem}
\begin{proof}
  \leanok
  Each \texttt{Tdc} is \texttt{TdcOf} (\ref{lem:m4-relations-aux}(1)); then \texttt{hcomm}.
\end{proof}

\begin{corollary}[The Hecke algebra on $H^n$ is commutative]\label{obl:algebra-model}
  % states [C] v1.18.0 obl:algebra (the ring-homomorphism consequence: the T's commute)
  \lean{GL2.M4.heckeAlg_comm_stmt, GL2.M4.heckeAlg_comm}
  \leanok
  \uses{thm:relations-comm-model, lem:m4-relations-aux}
  For $6\in R^\times$, \texttt{heckeAlg}~$\rho$~$n$ is commutative.
  \gtmodlit{\ref{lit:hecke-commute} as the hypothesis \gtlithyp{hcomm} (\texttt{HeckeCommute})}
\end{corollary}
\begin{proof}
  \leanok
  The generators commute pairwise (\ref{thm:relations-comm-model}); \ref{lem:m4-relations-aux}(5).
\end{proof}

\begin{definition}[The objects of \texttt{thm:classical} on the model]\label{def:m4-classical}
  % states [C] v1.18.0 thm:classical's objects: res, (conj_g)_*, the identification of the model's H^n with H^n(SL(2,Z);V)
  \lean{GL2.M4.VL, GL2.M4.Θ, GL2.M4.Spp, GL2.M4.mem_Spp, GL2.M4.conjHom, GL2.M4.conjCoeff, GL2.M4.resTo, GL2.M4.conjMap}
  \leanok
  \uses{def:m4-hecke, thm:model7-rmkorbifold}
  $V_L$ is $V$ as a representation of $\mathrm{SL}_2(\mathbb Z)$ through $\rho\circ e^{-1}$; $\Theta\colon H^n\to
  H^n(\mathrm{SL}_2(\mathbb Z);V_L)$ is $\Psi$ on cohomology followed by \ref{thm:model7-rmkorbifold} ($6\in R^\times$).
  $\Gamma''=\texttt{Spp}~g=\Gamma\cap g\Gamma g^{-1}$ (the group of the coset datum, $\Gamma=\bigsqcup\Gamma''s_j$) and
  $\Gamma'=\texttt{Spp}~g^{-1}=\Gamma\cap g^{-1}\Gamma g$; $(c_g)_*\colon H^n(\Gamma';V_L)\to H^n(\Gamma'';V_L)$ is mathlib's
  \texttt{groupCohomology.map} along $x\mapsto g^{-1}xg\colon\Gamma''\to\Gamma'$ with the coefficient map $\rho(g)$, which intertwines
  because $\rho(g)\rho(g^{-1}xg)=\rho(x)\rho(g)$; $\mathrm{res}$ is \texttt{groupCohomology.map} along the inclusion. In degree $0$,
  $\mathrm{cor}\circ(c_g)_*\circ\mathrm{res}$ is $v\mapsto\sum_j s_j^{-1}g\,v=\sum_j\rho(\delta_j)v$, the model's transport.
\end{definition}
\begin{proof}
  \leanok
  Definitions; \texttt{conjHom} and \texttt{conjCoeff} carry their membership and intertwining proofs.
\end{proof}

\begin{theorem}[$T_{\Gamma g\Gamma}$ has the shape $\mathrm{cor}\circ(c_g)_*\circ\mathrm{res}$ (weak form)]\label{thm:classical-model}
  % states a WEAK form of [C] v1.18.0 thm:classical on the model, with cor hypothesised (L20-0 §3.3; F3-20 MAJOR 1)
  % slot 23 (M6-A): proved outright, hcl is a theorem (thm:m6-hcl); GL2.M4.thm_classical (over hcl) is no longer linked
  \lean{GL2.M4.thm_classical_stmt, GL2.M6.thm_classical_outright}
  \leanok
  \uses{def:m4-classical, thm:m6-hcl, lem:m4-relations-aux}
  For $6\in R^\times$, $g\in\Delta$, every coset datum and every $\tau$: for \emph{some} $R$-linear
  $\mathrm{cor}\colon H^n(\Gamma'';V_L)\to H^n(\mathrm{SL}_2(\mathbb Z);V_L)$ with $\mathrm{cor}\circ\mathrm{res}=[\Gamma:\Gamma'']$,
  $\Theta\circ T_{\Gamma g\Gamma}=\mathrm{cor}\circ(c_g)_*\circ\mathrm{res}\circ\Theta$. This does not identify $\mathrm{cor}$
  with the corestriction ($\mathrm{cor}\circ\mathrm{res}=[\Gamma:\Gamma'']$ fixes $\mathrm{cor}$ on $\mathrm{im}\,\mathrm{res}$
  only) and is weaker than [C] \texttt{thm:classical}, which stays [C]-only: no \texttt{\textbackslash LEAN} or
  \texttt{\textbackslash LEANlit} marker on [C] \texttt{thm:classical} or \texttt{eq:corconjres} (F3-20 MAJOR 1).
  \emph{Restatement} (L20-0 \S3.3): mathlib \texttt{db584cd} has no corestriction, so $\mathrm{cor}$ is a hypothesised transfer,
  characterised only by $\mathrm{cor}\circ\mathrm{res}=[\Gamma:\Gamma'']$; [C] names the corestriction itself.
  \emph{Slot 23 (M6-A):} proved outright, with no literature hypothesis: \ref{lit:hecke-classical} is a theorem
  (\ref{thm:m6-hcl}, $\mathrm{cor}=$ \emph{the} corestriction), and [C] \texttt{thm:classical} in [C]'s sense is
  \ref{thm:m6-classical-theta}.
\end{theorem}
\begin{proof}
  \leanok
  \texttt{hcl} for the choice-free \texttt{TdcOf}; every coset datum and $\tau$ by \ref{lem:m4-relations-aux}(1). Since slot 23
  \texttt{hcl} is \ref{thm:m6-hcl}, so the row is \texttt{proved}.
\end{proof}

\begin{remark}[G5 rows without a Lean statement yet]\label{rmk:m4-g5-headers}
  \begin{itemize}
    \item [C] \texttt{thm:classical}: \textbf{open in [C]'s sense} (the identification with \emph{the} corestriction, Brown III.9).
      Only the weak form is stated: \ref{thm:classical-model}, \textsf{proved\_mod\_lit} since slot 20 (L20-5, branch H) over
      \ref{lit:hecke-classical}, where $\mathrm{cor}$ is pinned only by $\mathrm{cor}\circ\mathrm{res}=[\Gamma:\Gamma'']$ (F3-20
      MAJOR 1). Closing it (M6 scope at the earliest): pin $\mathrm{cor}$ by the degree-$0$ formula plus compatibility with the
      connecting maps (the transfer as a map of $\delta$-functors), or construct it on the bar complex (branch P). \emph{Slot 22
      (M5):} the genuine corestriction in every degree and the classical operator $\texttt{Tcls}$ with it now exist
      (\ref{def:m5-corH}, \ref{def:m5-Tcls}; FLT's \texttt{heckeH1} / \texttt{heckeInv} in degrees $1$ / $0$, \ref{thm:m5-heckeH1},
      \ref{thm:m5-delta0}); the identification $\Theta\circ T_{\Gamma g\Gamma}=\texttt{Tcls}\circ\Theta$ is still open (\ref{rmk:m5-open}).
      \emph{Slot 23 (M6-A):} proved, every $n$ (\ref{thm:m6-classical-theta}), by resolution transport.
    \item [C] \texttt{obl:algebra}, the ring homomorphism $\mathcal H(\Delta,\Gamma)\to\mathrm{End}(H^n)$ itself: no Hecke ring
      $\mathcal H(\Delta,\Gamma)$ exists in mathlib at \texttt{db584cd}. \emph{Corrected in slot 20} (L20-H,
      \texttt{docs/implementation/slot20/L20-H\_flt\_harvest.md}; the earlier sentence said none existed ``in the FLT harvest'' either, which is
      false): Anthropic FLT (upstream \texttt{6e837e75}, Apache-2.0) has the abstract ring, \texttt{HeckePair.HeckeAlgebra}
      (bi-invariant functions with finitely many cosets, convolution, \texttt{Ring} / \texttt{Algebra} instances, \texttt{heckeIndicator}), with
      Gelfand commutativity (\texttt{HeckePair.IsGelfandInvolution}, \texttt{heckeMul\_comm}), and the cohomological operator
      \texttt{HeckeCohomology.heckeH1}, the cocycle-level $\mathrm{cor}\circ(c_g)_\ast\circ\mathrm{res}$ in degree~$1$ (with
      \texttt{heckeH1\_delta0}, the $\delta$-compatibility the first item names as the way to pin $\mathrm{cor}$). Shimura's $R(\Gamma,\Delta)$
      is the subring of \texttt{HeckeAlgebra} supported on $\Delta$. What is missing upstream is the transport to $\mathrm{End}(H^n)$, i.e.\ the
      ring homomorphism itself. The FLT modules build on our toolchain (L20-H stage~2: $107/107$ Anthropic modules on v4.33.0 against
      \texttt{db584cd}); vendoring the import closure (25 modules, 2\,958 lines, with \texttt{MANIFEST}, \texttt{NOTICE}) is approved by the owner's licensing ruling of
      2026-10-04 (byte-identical copies, \texttt{NOTICE} / \texttt{ATTRIBUTION} as Apache-2.0 requires; recorded in
      \texttt{docs/implementation/slot20/owner\_rulings\_1004.md}) for slot~21 (M5). Until then the rows above state the homomorphism's consequences (\ref{thm:relations-psq-model} … \ref{obl:algebra-model}), and
      defining the ring here is not M4 scope.
  \end{itemize}
\end{remark}

\section{G9: level \texorpdfstring{$N$}{N} by Shapiro (headers)}

\begin{proposition}[Level-$N$ cohomology on the model is $H^n(\Gamma_0(N);\cdot)$]\label{prop:levelN-shapiro-model}
  % states [C] v1.18.0 def:levelellmodel's last claim / prop:lhcompliant (4), with mathlib's Coind in place of eq:MH
  % proved slot 20, L20-3 (M4-g): lean/GL2/M4/LevelN.lean
  \lean{GL2.M4.levelN_shapiro_stmt, GL2.M4.levelN_shapiro}
  \leanok
  \uses{thm:model7-rmkorbifold}
  For $6\in k^\times$, every $N$, every $k$-representation $A$ of $\Gamma_0(N)$ and every $n$:
  $H^n(\mathrm{relCxA}\ \mathrm{model7}\ (\mathrm{Coind}_{\Gamma_0(N)}^{\Gamma}A))\cong H^n(\Gamma_0(N);A)$. (Route:
  \texttt{rmkOrbifold} at the model, then mathlib's \texttt{groupCohomology.coindIso}.)
\end{proposition}
\begin{proof}
  \leanok
  \texttt{model7\_rmkOrbifold} composed with \texttt{groupCohomology.coindIso} ($\mathrm{SL}_2(\mathbb Z)$ and $k$ in \texttt{Type}).
\end{proof}

\begin{definition}[The explicit induced fibre; its isomorphism with \texttt{Coind}]\label{def:levelellmodel-model}
  % states [C] v1.18.0 def:inducedmodule, prop:blockmatrices (eq:MH), def:levelellmodel, prop:lhcompliant (4) in [C]'s fibre
  % slot 20, L20-3 (M4-g)
  \lean{GL2.M4.RTrans, GL2.M4.RTrans.split, GL2.M4.RTrans.ofFiniteIndex, GL2.M4.rtGamma0, GL2.M4.indRep, GL2.M4.IndN,
    GL2.M4.indToCoind, GL2.M4.indN_iso_coind, GL2.M4.indN_iso_coind_apply, GL2.M4.levelN_shapiro_ind}
  \leanok
  \uses{prop:levelN-shapiro-model}
  For $H\le G$ and a table $w_1,\dots,w_m$ of $H\backslash G$ (\texttt{RTrans}: $x=h(x)\,w_{j(x)}$ uniquely; one exists for every
  $H$ of finite index, \texttt{ofFiniteIndex}, e.g.\ $\Gamma_0(N)$, $N\ge1$), \texttt{IndN} is $\mathrm{Ind}_H^GA$ on $A^m$ with
  [C] \texttt{eq:MH}: $(\gamma f)_j=\rho(h_j(\gamma))f_{j\cdot\gamma}$, $w_j\gamma=h_j(\gamma)w_{j\cdot\gamma}$; the cocycle rule is
  \texttt{RTrans.split}. \texttt{indN\_iso\_coind}: $\mathrm{IndN}\cong\mathrm{Coind}_H^G A$, $f\mapsto(x\mapsto\rho(h(x))f_{j(x)})$.
  \texttt{levelN\_shapiro\_ind}: for $6\in k^\times$ and $H\le\mathrm{SL}_2(\mathbb Z)$ with a table,
  $H^n(\mathrm{relCxA}\ \mathrm{model7}\ \mathrm{IndN})\cong H^n(H;A)$, [C] \texttt{prop:lhcompliant}(4) in [C]'s fibre.
  \emph{Restatement.} Any table, not [C]'s fixed one (\texttt{def:cosetrepsN}); the isomorphism holds for every table.
\end{definition}
\begin{proof}
  \leanok
  Uniqueness of $x=h(x)w_{j(x)}$ gives the cocycle rule and both inverse laws; then \texttt{model7\_rmkOrbifold},
  \texttt{groupCohomology.functor} on the isomorphism, \texttt{coindIso}.
\end{proof}

\begin{proposition}[Level-$N$ correspondences: a cochain map; $\mathrm{cor}\circ(\mathrm{conj}_g)_*\circ\mathrm{res}$]\label{prop:levelcorrespondence-model}
  % states [C] v1.18.0 prop:levelcorrespondence (2), (3) (the function form under Frobenius reciprocity), (4)(b),(c); eq:levelcorr
  % slot 20, L20-3 (M4-h)
  \lean{GL2.Cone.GChainCx.resTo, GL2.M4.Bpp, GL2.M4.App, GL2.M4.RelTrans, GL2.M4.CochA, GL2.M4.resCx, GL2.M4.conjCx,
    GL2.M4.corFun_mem, GL2.M4.corCx, GL2.M4.levelFun, GL2.M4.levelCorr, GL2.M4.levelCorr_apply, GL2.M4.levelCorr_eq_res_cor,
    GL2.M4.levelCorr_isChainMap, GL2.M4.levelCorr_H_eq_res_cor, GL2.Cone.TwistedChainDatum.resTo,
    GL2.Cone.TwistedChainDatum.ofCentral, GL2.M4.levelCorr_central}
  \leanok
  \uses{def:cosettransport, def:levelellmodel-model}
  Let $A,B\le\Gamma$, $g\in\Delta$, $r_t$ a table of $B''\backslash B$, $B''=B\cap gAg^{-1}$ (\texttt{Bpp}), and $\tau$ a $g$-twisted
  chain datum (the level-$1$ one). On $\mathrm{Hom}_A(C_\bullet,V)$ (\texttt{CochA}), which is $\mathrm{Hom}_\Gamma(C_\bullet,
  \mathrm{Ind}_A^\Gamma V)$ under the cochain-level Frobenius reciprocity, \texttt{levelCorr} is
  $\psi\mapsto\bigl(y\mapsto\sum_t\rho(r_t^{-1}g)\,\psi(\tau(r_ty))\bigr)$, [C] \texttt{eq:levelcorr} in the form of
  \texttt{prop:levelcorrespondence}(3). It is a cochain map over every $R$, with no hypothesis on $\ell$ (\texttt{levelCorr\_isChainMap}),
  and, at cochain level and on $H^q$, $[A\,g\,B]=\mathrm{cor}^{B}_{B''}\circ(\mathrm{conj}_g)_*\circ\mathrm{res}^{A}_{A''}$,
  $A''=g^{-1}B''g$, with $\mathrm{res}$ the inclusion of equivariant cochains, $(\mathrm{conj}_g)_*\psi=\rho(g)\,\psi\circ\tau$ and
  $\mathrm{cor}\,\psi=\sum_t\rho(r_t^{-1})\,\psi(r_t\,\cdot)$, the instances (4)(b),(c) (\texttt{levelCorr\_eq\_res\_cor},
  \texttt{levelCorr\_H\_eq\_res\_cor}). For central $g$ with $\rho(g)=c$ and $\tau=\mathrm{id}$, $[A\,g\,A]=c$ (\texttt{levelCorr\_central}).
  \emph{Restatements.} (i) The block form \texttt{eq:levelcorr} on $\mathrm{Ind}_A^\Gamma V$ is not formalised separately: the row
  states (3)'s reciprocity form, which is the form \texttt{prop:levelcorrespondence}'s proof works in. (ii) $\tau$ over $\mathbb Z$
  (\texttt{TwistedChainDatum}), not over $R$. (iii) That the cochain-level $\mathrm{cor}$ and $(\mathrm{conj}_g)_*$ induce the
  group-cohomological ones on $H^\bullet(A;V)$ (Brown III.9--10) is not claimed here; it is the M4-f literature row.
\end{proposition}
\begin{proof}
  \leanok
  $B$-equivariance of $\mathrm{cor}$ is the computation of \texttt{lem:cosettransport}(1) at $g=1$ (reindex by $r_tb=b_tr_{f(t)}$);
  $(\mathrm{conj}_g)_*$ uses (T1); the factorisation is $\rho(r_t^{-1})\rho(g)=\rho(r_t^{-1}g)$; each factor commutes with $\partial$.
\end{proof}

\begin{definition}[Nebentypus data on the model]\label{def:nebdata-model}
  % states [C] v1.18.0 def:nebdata (2), (3), (5); lem:nebdata (1) (a d = 1 mod N on Gamma_0(N)), (2), (4)
  % slot 20, L20-3 (M4-h)
  \lean{GL2.M4.Δ0Q, GL2.M4.Δ0Q_le_ΔQ, GL2.M4.aZ, GL2.M4.aZ_eq, GL2.M4.aMod, GL2.M4.ρχ, GL2.M4.ρχ_apply, GL2.M4.Γ0Q,
    GL2.M4.Γ0Q_le_ΓQ, GL2.M4.Γ0Q_le_Δ0Q, GL2.M4.nebData}
  \leanok
  \uses{def:m4-hecke, conv:nebentypus}
  $\Delta_0(N)$ (\texttt{Δ0Q}): integral $\delta$ with $N\mid c_\delta$, $\gcd(a_\delta,N)=1$, a monoid; $\delta\mapsto a_\delta\bmod N$ is a
  monoid map (\texttt{aMod}); for a Dirichlet character $\chi$ (mathlib's \texttt{DirichletCharacter}~$R$~$N$),
  $\rho_\chi(\delta)=\chi(a_\delta)\rho(\delta)$ is a representation of $\Delta_0(N)$ (\texttt{ρχ}); $\Gamma_0(N)\subseteq\Delta_0(N)$,
  $a_\gamma d_\gamma\equiv1\ (N)$ on $\Gamma_0(N)$, and $\rho_\chi=\rho$ where $a_\delta\equiv1\ (N)$ (\texttt{nebData}).
  \emph{Restatement.} No $\det\delta>0$ (as for \texttt{ΔQ}); $\chi$ is any \texttt{DirichletCharacter}~$R$~$N$ over any commutative $R$, which
  contains [C]'s $\mathbb F_q$ and $W_K(\mathbb F_q)$ cases.
\end{definition}
\begin{proof}
  \leanok
  Entrywise congruences mod $N$; $\det=1$ gives $a d\equiv1$; $\chi$ and $\rho$ multiplicative.
\end{proof}

\begin{proposition}[Level-$N$ correspondences with a character]\label{prop:nebcorr-model}
  % states [C] v1.18.0 prop:nebcorr (1), (2) (cochain map, factorisation), (3) T(p,p) = chi(p) p^n, (4) term-wise; eq:nebcorr
  % slot 20, L20-3 (M4-h)
  \lean{GL2.M4.YN, GL2.M4.nebCorr, GL2.M4.nebCorr_isChainMap, GL2.M4.nebCorr_apply, GL2.M4.nebCorr_eq_res_cor,
    GL2.M4.nebCorr_apply_of_a_one, GL2.M4.nebTpp, GL2.M4.gP_mem_Δ0Q, GL2.M4.HN, GL2.M4.nebOp, GL2.M4.nebTprim}
  \leanok
  \uses{prop:levelcorrespondence-model, def:nebdata-model}
  For $A,B\le\Gamma_0(N)$, $g\in\Delta_0(N)$ and the level-$1$ $\tau$ (no new $\tau$-table), \texttt{nebCorr} is \texttt{levelCorr} at
  $\rho_\chi$: each term carries the scalar $\chi(a_{\delta_t})$, $\delta_t=r_t^{-1}g\in\Delta_0(N)$ (\texttt{nebCorr\_apply},
  [C] \texttt{eq:nebcorr}); a cochain map over every $R$; $=\mathrm{cor}\circ(\mathrm{conj}_g)_*\circ\mathrm{res}$; where every
  $a_{\delta_t}\equiv1$ the character drops out ((4), term-wise); a central $g$ with $\rho(g)=c$ acts by $\chi(a_g)c$, so
  $T^{(N,\chi)}(p,p)=\chi(p)p^n$ (\texttt{nebTpp}). \texttt{nebOp}, \texttt{nebTprim}: the induced operators on $H^n$ of the
  character model, $T^{(N,\chi)}_{p}$ / $U^{(N,\chi)}_p$ at $g=\mathrm{diag}(1,p^\mu)$.
  \emph{Restatement.} (4) is stated term-wise (the hypothesis $a_{\delta_t}\equiv1$), not as the restriction from the $\Gamma_1(N)$
  model, which is not formalised. \texttt{nebOp}, \texttt{nebTprim} are for a fixed table and $\tau$; independence of both (the
  $\Gamma_0(N)$ analogue of \texttt{thm:welldef}, [C] \texttt{prop:nebcorr}(2)'s ``independent of $\tau$ and of every choice of
  representatives'') is not formalised (F3-20 MINOR 6; an M5 item).
\end{proposition}
\begin{proof}
  \leanok
  Instances of \ref{prop:levelcorrespondence-model} on $\Gamma=\Gamma_0(N)$, $\Delta=\Delta_0(N)$, $\rho=\rho_\chi$.
\end{proof}

\begin{lemma}[Primitive products give the full recursion]\label{lem:m4-fullrec}
  % states the algebra of [C] v1.18.0 thm:relations (1)'s proof sketch (primitive -> full, with the (p+1) c term at mu = 0)
  % slot 20, L20-3
  \lean{GL2.M4.fullOf, GL2.M4.fullOf_add_two, GL2.M4.fullOf_rec}
  \leanok
  \uses{bridge:fullprimitive}
  In an $R$-algebra, with $P_0=1$, $F_\nu=\sum_{0\le2j\le\nu}c^jP_{\nu-2j}$: if $P_1^2=P_2+(p+1)c$ and
  $P_1P_{\mu+1}=P_{\mu+2}+pc\,P_\mu$ ($\mu\ge1$), then $F_{\nu+2}=P_1F_{\nu+1}-pc\,F_\nu$.
\end{lemma}
\begin{proof}
  \leanok
  $F_{\nu+2}=P_{\nu+2}+cF_\nu$; induction on $\nu$ in steps of two.
\end{proof}

\begin{theorem}[Hecke relations with a character]\label{thm:nebrelations-model}
  % states [C] v1.18.0 thm:nebrelations (2), (3); (1) only for T(p,p) (scalar, prop:nebcorr-model)
  % slot 20, L20-3 (M4-h)
  \lean{GL2.M4.nebRelations, GL2.M4.nebRelations_U}
  \leanok
  \uses{prop:nebcorr-model, lem:m4-fullrec}
  On $H^n$ of the character model, $6\in R^\times$: for $p\nmid N$ the full operators $T_{p^\nu}=\sum(\chi(p)c)^jT_{\Delta^{\mathrm{prim}}_{p^{\nu-2j}}}$
  satisfy $T_{p^{\nu+2}}=T_pT_{p^{\nu+1}}-\chi(p)\,p\,c\,T_{p^\nu}$ ($c=p^n$ for $\mathrm{Sym}^n$); for $p\mid N$, $U_{p^\nu}=U_p^\nu$.
  \gtmodlit{The double-coset products of the primitive operators with character in $\mathcal H(\Gamma_0(N),\Delta_0(N))$ (Shimura
  \S3.3 (3.3.2), (3.3.3), Thm.~3.24(5) with Prop.~3.17 via Prop.~3.31, Thm.~3.34(3) with Prop.~3.33; verified in R20-L), transported to $H^n$ (Brown III.9 with [C] \texttt{obl:algebra}), taken as
  the hypotheses \gtlithyp{hTpchi0} ($\mu=0$: $T_p^2=T_{\Delta^{\mathrm{prim}}_{p^2}}+(p+1)\chi(p)c$), \gtlithyp{hTpchi} ($\mu\ge1$) and
  \gtlithyp{hUchi} ($p\mid N$: $U_pU_{p^\mu}=U_{p^{\mu+1}}$).}
  \emph{Restatement.} (1), commutation of $T_p$ and $U_q$, is not stated (it would restate a hypothesis); $T(p,p)$ commutes with
  everything, being the scalar $\chi(p)c$ (\texttt{nebTpp}). The operators are those of a fixed table and $\tau$; independence of
  both (the $\Gamma_0(N)$ analogue of \texttt{thm:welldef}) is not formalised (F3-20 MINOR 6). The three literature identities
  are registered, with their R20-L-verified locators, in Part~L's slot-20 paragraph (F3-20 MINOR 9).
\end{theorem}
\begin{proof}
  \leanok
  \ref{lem:m4-fullrec} at $c\mapsto\chi(p)c$; induction on $\nu$ for $U$.
\end{proof}

\section{G9: Merel's set and the Heilbronn--Merel operator (M4-i, the L3-a oracle)}

\begin{definition}[Merel's set, executable]\label{def:merelset-model}
  % states [C] v1.18.0 lem:merelpath's set Mer_n, in the C++ order of hecke::level::MerelSet
  % slot 20, L20-4 (M4-i)
  \lean{GL2.M4.M2Z, GL2.M4.IsMerel, GL2.M4.toM2, GL2.M4.merelList, GL2.M4.mem_merelList, GL2.M4.merelSet,
    GL2.M4.mem_merelSet, GL2.M4.merelList_nodup, GL2.M4.merelList_sum}
  \leanok
  $\mathrm{Mer}_n=\{\smat abcd\in\mathrm M_2(\mathbb Z):ad-bc=n,\ a>b\ge0,\ d>c\ge0\}$, enumerated as \texttt{merelList}~$n$ ($a,d,b,c$
  ascending, the order of the C++ \texttt{MerelSet}); \texttt{mem\_merelSet}: membership is exactly Merel's conditions (the bound
  $a+d-1\le n$ is proved, not assumed); \texttt{merelList\_nodup}: no repetitions, so a sum over the list is the sum over the finite set.
\end{definition}
\begin{proof}
  \leanok
  $bc\le(a-1)(d-1)$ gives $n\ge a+d-1$; the four nested lists are disjoint by their first, last, second entries.
\end{proof}

\begin{lemma}[Merel's set and its path condition]\label{lem:merelpath-model}
  % states [C] v1.18.0 lem:merelpath (the path condition; the bound |Mer_n cap eta Gamma| <= x_0 is not stated)
  % slot 20, L20-4 (M4-i)
  \lean{GL2.M4.merel_path_stmt, GL2.M4.merel_path, GL2.M4.cuspOf, GL2.M4.epsD, GL2.M4.D0, GL2.M4.InCoset,
    GL2.M4.merelCoset, GL2.M4.succM, GL2.M4.predM, GL2.M4.succM_merel, GL2.M4.predM_merel, GL2.M4.pred_succ,
    GL2.M4.succ_pred, GL2.M4.upper_exists, GL2.M4.lower_exists, GL2.M4.upper_unique, GL2.M4.lower_unique}
  \leanok
  \uses{def:merelset-model}
  For $n\ge1$ and every $\eta\in\mathrm M_2(\mathbb Z)$ with $\det\eta=n$: $\sum_{\epsilon\in\mathrm{Mer}_n\cap\eta\Gamma}\epsilon\cdot D_0=D_0$ in the
  free abelian group on the cusps $\mathbb P^1(\mathbb Q)$ (\texttt{Option ℚ}, \texttt{none}${}=\infty$), $D_0=[\infty]-[0]$,
  $\epsilon\cdot D_0=[\epsilon\infty]-[\epsilon0]$ (tail minus head).
  \emph{Restatement.} [C]'s bound $\abs{\mathrm{Mer}_n\cap\eta\Gamma}\le x_0(\eta\mathbb Z^2)$ is not stated (nothing downstream uses it).
\end{lemma}
\begin{proof}
  \leanok
  Not [C]'s chain, its bookkeeping in one step. [C]'s successor (ii) $\epsilon\mapsto\epsilon\smat0{-1}1m$, $m=\lceil a/b\rceil$, maps the
  elements with $b>0$ into those with $c>0$, [C]'s predecessor (iii) $\epsilon\mapsto\epsilon\smat k1{-1}0$, $k=\lceil d/c\rceil$, maps
  back, and the two are mutually inverse (\texttt{pred\_succ}, \texttt{succ\_pred}: $\lceil(md-c)/d\rceil=m$ as $0\le c<d$); the
  successor's tail is the head. So the tails of the $c>0$ elements cancel the heads of the $b>0$ elements. What is left is the tail of
  the one element with $c=0$ (the upper Hermite form, $[\infty]$) and the head of the one element with $b=0$ (the lower Hermite form,
  $[0]$): each exists by Bézout (\texttt{upper\_exists}, \texttt{lower\_exists}) and is unique in its coset (two of them differ by
  $\delta\in\mathrm{SL}_2(\mathbb Z)$, which is forced to be $1$; \texttt{upper\_unique}, \texttt{lower\_unique}).
\end{proof}

\begin{proposition}[The Heilbronn--Merel operator over $\mathbb Z/q$, an executable spec]\label{prop:heilbronn-model}
  % states [C] v1.18.0 prop:heilbronn's definition eq:heilbronn and (4), and lem:cosetrepsN (4) for its terms; not (1), (2)
  % slot 20, L20-4 (M4-i)
  \lean{GL2.M4.heilbronn_spec_stmt, GL2.M4.heilbronn_spec, GL2.M4.Q4, GL2.M4.qmul, GL2.M4.qadj, GL2.M4.qdet,
    GL2.M4.rowMul, GL2.M4.primMod, GL2.M4.equivMod, GL2.M4.LevelTable, GL2.M4.clsOf, GL2.M4.termOf,
    GL2.M4.heilbronnTerms, GL2.M4.rhoZ, GL2.M4.heilbronnApply, GL2.M4.heilbronnSummand, GL2.M4.heilbronnApply_eq_sum}
  \leanok
  \uses{def:merelset-model, orc:rhoexec}
  For a table $w_1,\dots,w_m$ (\texttt{LevelTable}; $j\cdot\epsilon$ is the first $i$ with $(c_i,d_i)\equiv u\,(c_j,d_j)\epsilon\bmod N$,
  $u$ a unit), a prime-or-not $p$ and any $q$: \texttt{heilbronnApply}, $(Hx)_i=\sum\rho_n(h)x_j$ over the terms $(i,j,h)$ of
  \texttt{heilbronnTerms}, is the double sum of [C] \texttt{eq:heilbronn} over $\mathrm{Mer}_p\times\{0,\dots,m-1\}$ (the admissible pairs:
  $(c_j,d_j)\epsilon$ primitive modulo $N$); and on a table of $\mathrm{SL}_2(\mathbb Z)$ matrices every term matrix is
  $h=w_{j\cdot\epsilon}(\operatorname{adj}\epsilon)\operatorname{adj}(w_j)$ ($=w_{j\cdot\epsilon}(\operatorname{adj}\epsilon)w_j^{-1}$), of determinant
  $p$, with lower-left entry $\equiv0\bmod N$, and its index is $j\cdot\epsilon$.
  \emph{Restatement.} [C] (1) $H_{N,p}(\mathrm{Rel}_N)\subseteq\mathrm{Rel}_N$ and (2) $[H_{N,p}]=[A\,g\,A]$ on $M_N/\mathrm{Rel}_N$ are not
  stated: they need [C] \texttt{lem:manin} and \texttt{prop:levelorbit} (the cokernel presentation), which stay headers
  (\ref{rmk:m4-g9-headers}); the C++ checks them at its cells (\texttt{HeckeLevelNHeilbronnValidation}, (LN3)).
\end{proposition}
\begin{proof}
  \leanok
  The double sum: the term list is a \texttt{flatMap} over \texttt{merelList} of \texttt{filterMap}s, and \texttt{merelList} has no
  repetitions. The lower-left entry: with $(c_i,d_i)=u\,(c_j,d_j)\epsilon+N(k_1,k_2)$ it is $N(k_1X+k_2Y)$, because
  $(c_j,d_j)\,\epsilon\,(\operatorname{adj}\epsilon)\,(d_j,-c_j)^{\top}=p\,(c_jd_j-d_jc_j)=0$. The determinant: $1\cdot p\cdot1$.
\end{proof}

\begin{theorem}[The Heilbronn--Merel checker, green on real C++ output]\label{chk:merel-n5}
  % slot 20, L20-4 (M4-i): [LP] L3-a / L3-b on one HECKE_DUMP_MEREL dump of HeckeLevelNValidation
  \lean{GL2.M4.MerelDump, GL2.M4.tableCheck, GL2.M4.blockCheck, GL2.M4.denseOf, GL2.M4.merelCheck,
    GL2.M4.merelCheck_sound, GL2.M4.MerelDump.Sound, GL2.M4.fmod, GL2.M4.resNat, GL2.M4.Data.merel_N5_n2_p2_check, GL2.M4.Data.merel_N5_n2_p5_check,
    GL2.M4.Data.merel_N5_n2_p2_sound, GL2.M4.Data.merel_N5_n2_p5_sound, GL2.M4.Data.merel_tampered_rejected}
  \leanok
  \uses{prop:heilbronn-model, lem:merelpath-model, chk:rhoexec}
  The C++ \texttt{HeilbronnOperator} at \texttt{HeckeLevelNValidation}'s own cell $(N,n)=(5,2)$ over $\mathbb F_{1000003}$, for $p=2$
  ($T_2^{(5)}$) and $p=5$ ($U_5^{(5)}$, inadmissible pairs dropped), dumped by \texttt{HECKE\_DUMP\_MEREL}: the table is a transversal
  (in $\mathrm{SL}_2(\mathbb Z)$, pairwise inequivalent bottom rows, every primitive class hit), $\mathrm{Mer}_p$ and the term list are
  the spec's, every block is $\rho_n(h\bmod q)$, and the dense $D\times D$ matrix is, entrywise in $\mathbb Z/q$, the matrix of
  \texttt{heilbronnApply} on the dumped table (\texttt{merelCheck\_sound}). One block entry moved by one is rejected.
  \emph{Provenance} (F3-20 MINOR 7): the C++ \texttt{MEREL\_DUMP} line prints \texttt{dense\_fnv64}, the FNV-1a-64
  (\texttt{Fnv.h}) of $D$ and the dense entries row-major, and \texttt{merel\_dump\_to\_lean.py} recomputes it into the generated
  docstrings: \texttt{0x029f2ae7b5c322c1} ($p=2$) and \texttt{0xa321b63a9401bf42} ($p=5$), equal on both sides (slot-20 L20-6
  repairs; the dump files are byte-identical to L20-4's).
\end{theorem}
\begin{proof}
  \leanok
  \texttt{decide +kernel} on the generated data \texttt{GL2/Checkers/MerelData.lean} (\texttt{lean/tools/merel\_dump\_to\_lean.py}); then
  \texttt{merelCheck\_sound}.
\end{proof}

\begin{remark}[G9 rows without a Lean statement yet (after slot 20, L20-3 and L20-4)]\label{rmk:m4-g9-headers}
  Each of the following needs a definition first. The task list (\texttt{L19-5\_M4\_skeleton.md} §3) gives the Lean names.
  \begin{itemize}
    \item [C] \texttt{def:levelellmodel}, \texttt{prop:levelcorrespondence}: rows \ref{def:levelellmodel-model},
      \ref{prop:levelcorrespondence-model} (slot 20, L20-3).
    \item [C] \texttt{prop:levelorbit}: orbit-wise fixed points and $H^1$ as a cokernel.
    \item [C] \texttt{lem:manin}: Manin's sequence. [C] \texttt{lem:merelpath} and \texttt{prop:heilbronn} (its definition
      [C] \texttt{eq:heilbronn}, (4) and the terms' \texttt{lem:cosetrepsN}(4), not (1), (2)) are rows \ref{def:merelset-model},
      \ref{lem:merelpath-model}, \ref{prop:heilbronn-model}, with the [LP] L3-a checker \ref{chk:merel-n5} (slot 20, L20-4);
      \texttt{prop:heilbronn}(1), (2) wait for \texttt{lem:manin} and \texttt{prop:levelorbit}.
    \item [C] \texttt{sec:nebentypus}: \texttt{def:nebdata}, \texttt{prop:nebcorr}, \texttt{thm:nebrelations} are rows
      \ref{def:nebdata-model}, \ref{prop:nebcorr-model}, \ref{thm:nebrelations-model} (slot 20, L20-3); \texttt{prop:neborbit}
      and the $\Gamma_1(N)$ model (\texttt{lem:gammaonetable}, \texttt{prop:nebdecomp}) stay headers. They share the conventions of
      \ref{conv:nebentypus} and \ref{bridge:nebentypus-dictionary}.
  \end{itemize}
\end{remark}
